\section{Solution Process}

After the variable system and candidate set are fixed, the solution process should be presented in the same order as the decision burden actually unfolds. The first computational layer produces the supplier scores and ranking order. Its purpose is to compress the comparison information into a manageable candidate set, not to finish the paper. Accordingly, the solution description should state how the score is computed, how the top candidates are retained, and what threshold or rule is used to enter the planning stage.

The second computational layer solves the executable allocation problem on the reduced candidate set. At this point the paper should explain how demand, capacity, and effective loss constraints are inserted into the model, how the objective reflects total burden or operating cost, and which solver or heuristic is used to obtain the plan. Even when the final implementation is concise, the solution section should still make the model-to-output path visible: candidate set in, feasible plan out.

The third computational layer is the validation or comparison layer. Once the baseline plan is obtained, the same pipeline is rerun under disturbed demand, altered loss rates, or different capacity assumptions. The goal of this step is not to create extra pages, but to show whether the recommended supplier structure and allocation pattern remain stable when the environment changes within a reasonable range. In a serious team-written paper, this layer is explicitly presented as part of the solution route rather than as an afterthought.

From a reproducibility perspective, the implementation should generate stable artifacts for each stage: a ranking table for the evaluation layer, a decision table for the executable plan, and a scenario or feasibility artifact for validation. When these outputs are generated by one entry point and recorded in the artifact ledger, the paper gains a much stronger research value, because the numerical claims in the abstract and conclusion can be traced back to concrete files instead of staying at the level of prose.
